\documentclass[12pt, a4paper]{article}

% ---- Packages ----
\usepackage[T1]{fontenc}
\usepackage[utf8]{inputenc}
\usepackage{amsmath, amssymb}
\usepackage{booktabs}
\usepackage{geometry}
\usepackage{hyperref}
\usepackage{xcolor}
\usepackage{array}
\usepackage{caption}
\usepackage{longtable}

\geometry{margin=2.5cm}
\hypersetup{colorlinks=true, linkcolor=blue, citecolor=blue, urlcolor=blue}

% ---- Title ----
\title{\textbf{Phase Equilibrium Concentration Derivation\\for Binary Alloy Systems}\\[0.5em]
\large Notes on Thermodynamics and Materials Science}
\author{}
\date{}

\begin{document}

\maketitle
\tableofcontents
\newpage

%% ==============================================================
\section{Known Conditions and Symbol Definitions}
%% ==============================================================

Consider a binary alloy system containing two solute elements (denoted as elem~1 and elem~2).
The physical parameters are defined as follows.

\subsection{Initial State}
\begin{itemize}
  \item Initial concentrations: $C_1^0$ (elem~1), $C_2^0$ (elem~2)
  \item Initial temperature: $T_0$ (liquidus temperature of the base matrix)
\end{itemize}

\subsection{Physical Parameters}
\begin{itemize}
  \item Liquidus slopes: $m_1$ (elem~1), $m_2$ (elem~2) \quad [K/at.\%]
  \item Equilibrium partition coefficients: $k_1$ (elem~1), $k_2$ (elem~2) \quad [dimensionless]
  \item Temperature correction term: $\Delta T_R$ (known constant accounting for matrix liquidus shift)
\end{itemize}

\subsection{Current State (at a given temperature $T$)}
\begin{itemize}
  \item Actual liquid concentrations: $C_1$, $C_2$
\end{itemize}

\subsection{Quantities to Solve}
\begin{itemize}
  \item Equilibrium interface concentrations: $C_1^{eq}$, $C_2^{eq}$
  \item Solid fraction increment: $\Delta f_S$
\end{itemize}

%% ==============================================================
\section{Governing Equations}
%% ==============================================================

Two core equations are established based on the linear thermodynamic approximation
and the solute redistribution principle.

\subsection{Equation~1: Temperature--Concentration Linearization}

The local liquidus temperature is assumed to shift linearly with solute concentration,
with contributions from each element superimposed:
\begin{equation}
  T = T_0
    + m_1\!\left(C_1^{eq} - C_1^0\right)
    + m_2\!\left(C_2^{eq} - C_2^0\right)
    + \Delta T_R .
  \label{eq:temp}
\end{equation}

\subsection{Equation~2: Solid Fraction from Partition Coefficients}

Based on the lever rule and the definition of the partition coefficient,
the solid fraction increment $\Delta f_S$ must be consistent for both elements:
\begin{equation}
  \Delta f_S
  = \frac{C_1^{eq} - C_1}{C_1^{eq}(k_1 - 1)}
  = \frac{C_2^{eq} - C_2}{C_2^{eq}(k_2 - 1)} .
  \label{eq:fs}
\end{equation}

%% ==============================================================
\section{Derivation}
%% ==============================================================

\subsection{Step 1: Relating $C_1^{eq}$ and $C_2^{eq}$ via Equation~\eqref{eq:fs}}

Cross-multiplying Eq.~\eqref{eq:fs}:
\[
  C_2^{eq}(k_2 - 1)(C_1^{eq} - C_1)
  = C_1^{eq}(k_1 - 1)(C_2^{eq} - C_2).
\]

Let $A_i = k_i - 1$ for brevity. Expanding and rearranging yields an expression
for $C_2^{eq}$ as a function of $C_1^{eq}$:
\begin{equation}
  C_2^{eq}
  = \frac{A_1\, C_2\, C_1^{eq}}
         {A_1\, C_1^{eq} - A_2\!\left(C_1^{eq} - C_1\right)} .
  \label{eq:C2eq}
\end{equation}

Substituting $A_1 = k_1 - 1$ and $A_2 = k_2 - 1$:
\begin{equation}
  C_2^{eq}
  = \frac{(k_1-1)\,C_2\,C_1^{eq}}
         {(k_2-1)\,C_1^{eq} + (k_1 - k_2)\,C_1} .
  \label{eq:C2eq_explicit}
\end{equation}

\subsection{Step 2: Substituting into Equation~\eqref{eq:temp}}

Define the effective undercooling
\[
  \Delta T \;=\; T - T_0 - m_1 C_1^0 - m_2 C_2^0 - \Delta T_R .
\]

Substituting Eq.~\eqref{eq:C2eq_explicit} into Eq.~\eqref{eq:temp} and letting
the denominator $D = (k_2-1)\,C_1^{eq} + (k_1-k_2)\,C_1$, after multiplying through
by $D$ and collecting terms in $C_1^{eq}$, one obtains the quadratic equation:
\begin{equation}
  a\,\bigl(C_1^{eq}\bigr)^2 + b\,C_1^{eq} + c = 0 .
  \label{eq:quadratic}
\end{equation}

%% ==============================================================
\section{Closed-Form Results}
%% ==============================================================

\subsection{Quadratic Coefficients}

\begin{align}
  a &= m_1(k_2-1) + m_2(k_1-1)\,\frac{C_2}{C_1} , \label{eq:coef_a} \\[4pt]
  b &= m_1(k_1-k_2)\,C_1 + m_2(k_1-1)\,C_2 - \Delta T\,(k_2-1) , \label{eq:coef_b} \\[4pt]
  c &= -\Delta T\,(k_1-k_2)\,C_1 , \label{eq:coef_c}
\end{align}
where
\[
  \Delta T = T - T_0 - m_1 C_1^0 - m_2 C_2^0 - \Delta T_R .
\]

\subsection{Solution for $C_1^{eq}$}

\begin{equation}
  \boxed{C_1^{eq} = \frac{-b + \sqrt{b^2 - 4ac}}{2a}}
  \label{eq:C1eq}
\end{equation}

\textbf{Remark.} The positive root is selected because concentrations must be positive
(typically $C_1^{eq} > C_1$ during solidification). If $a \to 0$ (i.e.\ $k_1 \approx k_2$),
the equation degenerates to a linear problem solved directly from Eq.~\eqref{eq:coef_b}.

\subsection{Solution for $C_2^{eq}$}

\begin{equation}
  \boxed{C_2^{eq} = \frac{(k_1-1)\,C_2\,C_1^{eq}}
                         {(k_2-1)\,C_1^{eq} + (k_1-k_2)\,C_1}}
  \label{eq:C2eq_final}
\end{equation}

\subsection{Solid Fraction Increment}

\begin{equation}
  \boxed{\Delta f_S = \frac{C_1^{eq} - C_1}{C_1^{eq}(k_1-1)}}
  \label{eq:DeltaFS}
\end{equation}

%% ==============================================================
\newpage
\section{Parameter Set: CoCrNi-Ti-C System}
%% ==============================================================

This section extracts parameters for the CoCrNi-Ti-C alloy system from the report:\\
\textit{``3D Cellular Automaton Simulation of Dendrite Growth in CoCrNi-Ti-C Alloys:
A Comprehensive Parameterization Framework''} (March 2026).\\
These parameters are used to populate $m_1$, $m_2$, $k_1$, $k_2$, and related
quantities in the binary derivation above.

\subsection{System Definition}

The CoCrNi-Ti-C system is treated as a pseudo-binary system:
\begin{itemize}
  \item elem~1 = Titanium (Ti)
  \item elem~2 = Carbon (C)
  \item Matrix = equiatomic CoCrNi (effective solvent)
\end{itemize}

\subsection{Thermodynamic Parameters}

\begin{table}[htbp]
\centering
\caption{Thermodynamic parameters ($m_i$, $k_i$) for the CoCrNi-Ti-C system.}
\label{tab:thermo}
\begin{tabular}{lllll}
\toprule
Symbol & Description & Ti (elem~1) & C (elem~2) & Source \\
\midrule
$k_i$  & Equilibrium partition coefficient & $0.06$--$0.12$ & $0.03$--$0.08$ & CALPHAD (TCHEA6); [9,15] \\
$m_i$  & Liquidus slope (K/at.\%)          & $-20$ to $-35$ & $-40$ to $-80$ & CALPHAD; [9,13,15] \\
$T_0$  & CoCrNi matrix liquidus temperature & \multicolumn{2}{c}{1610 K (1337 °C)} & Exp.\ DSC + CALPHAD [10] \\
$T_\text{sol}$ & Solidus temperature        & \multicolumn{2}{c}{1443 K (1170 °C)} & Exp.\ DSC + CALPHAD [10] \\
\bottomrule
\end{tabular}
\end{table}

\noindent\textbf{Notes:}
\begin{itemize}
  \item $k < 1$: both Ti and C strongly partition into the liquid phase.
  \item $m < 0$: solute enrichment depresses the liquidus temperature.
  \item Precise values should be extracted from Thermo-Calc + TCHEA6 at the exact composition.
\end{itemize}

\subsection{Temperature Correction Term $\Delta T_R$}

For a given alloy $(\text{CoCrNi})_{1-2x}\text{Ti}_x\text{C}_x$, the liquidus temperature
shifts relative to the pure CoCrNi matrix:
\[
  \Delta T_R = (m_1 + m_2)\,x .
\]
For the representative composition $x = 0.004$ (0.4 at.\% Ti and C each):
\[
  \Delta T_R \approx (-25 + (-60)) \times 0.4 = -34\ \text{K} .
\]

\subsection{Initial Concentrations}

\begin{table}[htbp]
\centering
\caption{Initial alloy concentrations.}
\label{tab:init}
\begin{tabular}{llll}
\toprule
Symbol & Description & Typical value & Unit \\
\midrule
$C_1^0$ & Ti initial concentration & 0.4 (range 0.2--1.0) & at.\% \\
$C_2^0$ & C initial concentration  & 0.4 (equimolar with Ti) & at.\% \\
\bottomrule
\end{tabular}
\end{table}

\subsection{Diffusion Kinetics Parameters}

\begin{table}[htbp]
\centering
\caption{Diffusion coefficients in liquid and solid phases.}
\label{tab:diff}
\begin{tabular}{llllll}
\toprule
Symbol & Description & Ti & C & Unit & Source \\
\midrule
$D_L$ & Liquid diffusivity & $3\times10^{-9}$ & $8\times10^{-9}$ & m$^2$/s & [13,16] \\
$D_S$ & Solid diffusivity  & $\approx 0$ (Scheil) & $10^{-14}$--$10^{-13}$ & m$^2$/s & DICTRA (MOBHEA2) [9] \\
\bottomrule
\end{tabular}
\end{table}

\noindent\textbf{Note.} Due to the sluggish diffusion effect in CoCrNi, solid back-diffusion
$D_S$ for Ti is safely neglected (Scheil--Gulliver approximation).

\subsection{Thermophysical Properties}

\begin{table}[htbp]
\centering
\caption{Additional thermophysical parameters.}
\label{tab:thermo2}
\begin{tabular}{lllll}
\toprule
Symbol & Parameter & Value & Unit & Source \\
\midrule
$L$          & Latent heat of fusion           & $2.35\times10^5$          & J/kg        & Exp.\ [10] \\
$\rho$       & Density                          & $8100$--$8300$            & kg/m$^3$    & [21] \\
$c_p$        & Specific heat capacity           & $530$--$560$              & J/(kg$\cdot$K) & [10] \\
$\sigma_{SL}$& Solid--liquid interfacial energy & $0.18$--$0.25$            & J/m$^2$     & Phase-field inversion [17] \\
$\Gamma$     & Gibbs--Thomson coefficient       & $(1\text{--}3)\times10^{-7}$ & K$\cdot$m  & [18] \\
$\varepsilon_4$ & Anisotropy strength           & $0.03$--$0.05$            & --          & [13] \\
\bottomrule
\end{tabular}
\end{table}

\subsection{Nucleation Parameters (for CET Modeling)}

\begin{table}[htbp]
\centering
\caption{Nucleation parameters for the columnar-to-equiaxed transition (CET) model.}
\label{tab:nucl}
\begin{tabular}{llll}
\toprule
Parameter & Description & Recommended value & Unit \\
\midrule
$n_{\max,v}$      & Volumetric nucleation density        & $10^{11}$--$10^{13}$ & m$^{-3}$ \\
$\Delta T_{N,v}$  & Mean volumetric nucleation undercooling & $1.5$--$4$          & K \\
$\sigma_{\Delta T,v}$ & Std.\ dev.\ of nucleation distribution & $0.5$--$1$      & K \\
$n_{\max,s}$      & Surface nucleation density           & $10^{12}$--$10^{14}$ & m$^{-2}$ \\
$\Delta T_{N,s}$  & Surface nucleation undercooling      & $\approx 0.5$        & K \\
\bottomrule
\end{tabular}
\end{table}

\subsection{Recommended Parameter Values for the Binary Derivation}

Substituting the above into the equations of Sections~2--4:
\begin{align*}
  m_1 &= -25\ \text{K/at.\%} \quad \text{(Ti liquidus slope, recommended)} \\
  m_2 &= -60\ \text{K/at.\%} \quad \text{(C liquidus slope, recommended)} \\
  k_1 &= 0.08 \quad \text{(Ti partition coefficient, recommended)} \\
  k_2 &= 0.05 \quad \text{(C partition coefficient, recommended)} \\
  T_0 &= 1610\ \text{K} \quad \text{(CoCrNi matrix liquidus)} \\
  \Delta T_R &= \text{from CALPHAD, or approximated as above}
\end{align*}

The solved quantities $C_1^{eq}$ and $C_2^{eq}$ represent the equilibrium liquid concentrations
of Ti and C at the solid--liquid interface; $\Delta f_S$ is the solid fraction increment
at each time step, reflecting the advancement rate of the solidification front.

\bigskip
\noindent\textit{Parameter source: ``3D CA Simulation of Dendrite Growth in CoCrNi-Ti-C Alloys,'' March 2026.}

\end{document}
